% Source: https://artofproblemsolving.com/community/c3h1438519_taking_crt_in_mods
% Stable identity: aops-topic-1438519-q01
\textbf{Modulo $8$:} $2^{192}\equiv0$.

\textbf{Modulo $125$:} by Euler's theorem $2^{100}\equiv1\pmod{125}$ (as $\varphi(125)=100$), so $2^{192}\equiv2^{92}$, and $2^{92}\cdot2^{8}=2^{100}\equiv1$. Now $2^7=128\equiv3$, so $2^8\equiv6$, and since $6\cdot21=126\equiv1$, we get
\[
2^{192}\equiv2^{92}\equiv6^{-1}\equiv21\pmod{125}.
\]
\textbf{Combine:} $x=21+125t$ with $x\equiv0\pmod8$. As $21\equiv5$ and $125\equiv5\pmod8$, we need $5+5t\equiv0$, so $t\equiv7\pmod8$, and $x=21+875=896$. The remainder is $\boxed{896}$.

\textbf{The student's method is wrong.} Knowing a number modulo $5$ does not determine it modulo $125$, and cubing does not turn a congruence mod $5$ into one mod $125$ (for example $6\equiv1\pmod5$ but $6\not\equiv1\pmod{125}$). Here $2^{192}\equiv1\pmod5$ is true, but $2^{192}\equiv21\pmod{125}$, so the answer $376$ is incorrect.
